\subsection{ANGULAR MEASURE}

By Theorem 1 of no. 1, the topological group \(\mathfrak{A}\) of angles is isomorphic to \(\mathbf{T}\) . Every strict morphism of \(\mathbf{R}\) onto \(\mathfrak{A}\) can be obtained by composing the isomorphism \(z \to \operatorname{Am}(z)\) of \(\mathbf{U}\) onto \(\mathfrak{A}\) with a strict morphism of \(\mathbf{R}\) onto \(\mathbf{U}\) ; if we put \(\Im(x) = \operatorname{Am}(e(x))\) , every strict morphism of \(\mathbf{R}\) onto \(\mathfrak{A}\) is therefore of the form \(x \to \Im(x/a)\) ( \(a \neq 0\) ). Given a real number \(a > 0\) , fixed once and for all, every angle \(\theta\) corresponds, by the homomorphism \(x \to \Im(x/a)\) , to a class of real numbers mod \(a\) (i.e., an element of \(\mathbf{Z}/a\mathbf{Z}\) ) which is called the measure of \(\theta\) relative to the base \(a\) ; by abuse of language, every real number in this class is also called a measure of \(\theta\) ; the angle \(\Im(x/a)\) is called the angle with measure \(x\) (relative to the base a). If x is a measure of \(\theta\) and \(x'\) a measure of \(\theta'\) (relative to the same base) then \(x + x'\) is a measure of \(\theta + \theta'\) , and -x is a measure of - \(\theta\) . The principal measure of an angle (relative to the base a) is that one of its measures which lies in the interval {[}0, a{]}.

Choice of a base \(a\) . We restrict ourselves always to bases \(a > 1\) . To each \(a > 1\) corresponds an angle \(\omega = \Im(1/a)\) whose principal measure is \(1\) , and which is called the unit of angular measure relative to the base \(a\) ; conversely, to each angle \(\omega \neq 0\) there corresponds a unique \(a > 1\) such that \(\Im(1/a) = \omega\) , so that knowledge of the unit of angular measure determines the base \(a > 1\) entirely.

In numerical calculations one usually takes either \(a = 360\) or \(a = 400\) ; the corresponding unit of angular measure is called the degree ( \(a = 360\) ) or the grade ( \(a = 400\) ).

In analysis, and indeed in all branches of mathematics where numerical calculation is not involved, the base a defined by the condition

\[
\lim _ {x \rightarrow 0} \frac {e (x / a) - 1}{x} = 1
\]

is universally used; this base is denoted by \(2\pi\) . The corresponding unit of angular measure is called a radian, and the measure is called radian measure; with the definition of \(e^{z}\) for complex z mentioned earlier, we have \(\boldsymbol{e}(x)=e^{2\pi ix}\) for all \(x\in R\) .

Once the base a has been chosen, when one speaks of an angle one usually means a measure of this angle relative to the base a; this abuse of language has no drawback provided (as is always the case when numerical calculations are not involved) the base a remains fixed throughout, and provided that one remembers that two real numbers which are congruent mod a correspond to the same angle.

For example, what is usually understood by the amplitude of a complex number \(z \neq 0\) is a radian measure of this angle, determined by conventions which will depend on the question under consideration; once these conventions have been made, the measure of the amplitude thus chosen is denoted by Am (z).

